\chapter{Metodología}\label{cap:metodologia}
% GUÍA — Qué va en este capítulo:
%  - La metodología de trabajo seguida (p. ej. un proceso iterativo o ágil)
%    y cómo se ha aplicado en la práctica: iteraciones, planificación, roles.
%  - Las herramientas, lenguajes y recursos hardware/software utilizados.
%  - Si procede, el diseño de la solución (arquitectura, modelos, diagramas).
%  - Este capítulo incluye ejemplos de figura, tabla, cita, listado y acrónimo
%    para que copies su estructura; bórralos cuando ya no los necesites.

Sustituye este párrafo por la descripción de cómo has llevado a cabo el
trabajo.

\section{Ejemplos para copiar}

% --- CITA: \cite{clave}, donde «clave» es la primera palabra de la entrada
%     en bibliografia.bib. El ~ evita que la cita quede sola en otra línea.
El ciclo de trabajo se inspira en los métodos iterativos descritos por
Sommerville~\cite{sommerville2016}.

% --- ACRÓNIMO: \ac{SIGLA}; la sigla debe estar definida en 0-inicio/acronimos.tex.
Los modelos del diseño se expresan en \ac{UML}.

% --- REFERENCIAS CRUZADAS: \ref{etiqueta} da el número de una figura, tabla,
%     listado o capítulo que tenga \label{etiqueta}.
La figura~\ref{fig:ciclo} resume el proceso, la tabla~\ref{tab:herramientas}
recoge las herramientas empleadas y el listado~\ref{lst:hola} muestra un
fragmento de código.

% --- FIGURA: el archivo va en la carpeta figuras/ (PDF, PNG o JPG) y se
%     escribe sin la carpeta. width=… fija el ancho (aquí, 80 % del texto).
\begin{figure}[htbp]
  \centering
  \includegraphics[width=0.8\textwidth]{ejemplo}
  \caption{Ciclo de trabajo iterativo seguido en el proyecto.}
  \label{fig:ciclo}
\end{figure}

% --- TABLA: columnas l (izquierda), c (centro), r (derecha); & separa
%     celdas y \\ termina la fila. \toprule, \midrule y \bottomrule son las
%     líneas horizontales.
\begin{table}[htbp]
  \centering
  \caption{Herramientas utilizadas.}
  \label{tab:herramientas}
  \begin{tabular}{lll}
    \toprule
    \textbf{Herramienta} & \textbf{Uso}           & \textbf{Versión} \\
    \midrule
    Python               & Desarrollo             & 3.12 \\
    Git                  & Control de versiones   & 2.45 \\
    \LaTeX{}             & Redacción de la memoria & TeX Live 2026 \\
    \bottomrule
  \end{tabular}
\end{table}

% --- LISTADO DE CÓDIGO: language=… activa el resaltado (Python, Java, C, C++,
%     SQL, HTML…). El código se copia tal cual entre begin y end.
\begin{lstlisting}[language=Python, caption={Ejemplo de función en Python.}, label={lst:hola}]
def saludar(nombre):
    """Devuelve un saludo para la persona indicada."""
    return f"Hola, {nombre}"
\end{lstlisting}
